%% kapa: exam=1 round=30 year=2019 period=1 form=A subject=real_estate subject_ko=부동산학원론 seq=14 number=94
%% source: Q-Net 감정평가사 기출문제 (cst00309) artlSeq=5205205 file=extracted/1008/5205205/2019년 제30회 감정평가사 1차시험 문제지/2019년 제30회 감정평가사 1차시험 문제지/2019년도 제30회 감정평가사 1차 1교시 A형.hwp
%% conversion: kordoc 4.24.51 -> tools/gen1.py (auto); review=proofread-sonnet; flags=none
\begin{question}{94}
부동산 증권에 관한 설명으로 옳지 않은 것은?
\begin{choices}
  \item 자산유동화증권(ABS)은 금융기관 및 기업이 보유하고 있는 매출채권, 부동산저당채권 등 현금흐름이 보장되는 자산을 담보로 발행하는 증권을 의미한다.
  \item 저당담보부채권(MBB)은 모기지풀에서 발생하는 현금흐름과 관련된 위험을 투자자에게 이전하는 채권이다.
  \item 주택저당증권(MBS)은 금융기관 등이 주택자금을 대출하고 취득한 주택저당채권을 유동화전문회사 등이 양수하여 이를 기초로 발행하는 증권을 의미한다.
  \item 저당이체증권(MPTS)은 발행기관이 원리금수취권과 주택저당권에 대한 지분권을 모두 투자자에게 이전하는 증권이다.
  \item 다계층증권(CMO)은 저당채권의 발행액을 몇 개의 계층으로 나눈 후 각 계층마다 상이한 이자율을 적용하고 원금이 지급되는 순서를 다르게 정할 수 있다.
\end{choices}
\end{question}
\answer{2}
